\section{Conclusion}\label{sec:conclusion}

This study shows that an LLM-based detector can identify and locate SQL antipatterns in the Java expressions and generated schema classes through which jOOQ represents database access. The selected detector achieved micro F1 \num{0.869} on held-out projects, with per-class F1 ranging from \num{0.481} to \num{0.974}. These results support using its detections as starting points for targeted review. Their interpretation depends on the antipattern class, the available source context, and the decisions encoded in the class definition. Inspection of disagreements also exposed missing reference annotations, showing why evaluating such a detector requires scrutiny of both its predictions and the human reference.

The corpus findings suggest different scopes for that review. ID Required flags covered almost half of generated-schema files and generally occurred once per flagged file. Implicit Columns flags covered about one quarter of application/query files and averaged \num{2.80} flags per flagged file. These patterns point towards reviewing schema decisions across tables and examining related query expressions together. Repository flag totals also reflected both code volume and flag density. The practical lesson is that deciding what to inspect requires knowing what was flagged and how the findings are distributed.

An independent audit of flagged and unflagged code is the next step towards establishing how these findings generalise across the corpus. The \texttt{selectFrom} and \texttt{select()} fragment categories accounted for \qty{71.7}{\percent} of Implicit Columns flags, providing a focused starting point for that audit. The observed detector errors identify additional source forms to examine for missed occurrences. A comparison with explicit rules on the same annotated code could then establish where an LLM contributes useful interpretation. The longer-term aim is to give developers findings whose meaning and limits are clear enough to support informed changes to database-access code.
